\section{Related work}

The closest causal comparison is the contaminated-treatment dose-response estimator of \citet{HuangZhang2023}. Under their positivity, smoothing, sieve, and supersmooth-error conditions, the fixed-dose expansion in arXiv version 2, Theorem 4.4 and equation (21) has bias orders
\[
h^2+h_0^2
\]
and stochastic remainder
\[
O_p\!\left\{\frac{N^{-1/2}}{\min(e(h),e(h_0))}\right\},
\qquad
e(b)=b^{1/2}\exp(-b^{-r}/\gamma),
\]
where \(h,h_0\) are bandwidths, \(N\) is sample size, \(r\) is the supersmooth exponent, and \(\gamma\) is the error-tail constant; the normal limit also uses the theorem's variance lower bound \citep[arXiv version 2, Theorem 4.4, equation (21)]{HuangZhang2023}. The present numerical bounds and schedule-level exchangeability specify a different structural class. Within that class, the delivered advance is a minimax characterization uniform over the known noise scale, allowing polynomially weak dose support and matching the absolute-risk order with honest connected-interval length.
% [Huang and Zhang, arXiv version 2, Theorem 4.4](https://arxiv.org/pdf/2211.04642v2)

The error-free endpoint connects to pointwise regression with a degenerate design. For sample size \leanref{sym:n}{\ensuremath{n}}, mean smoothness exponent \leanref{sym:beta(smoothness exponent in (0,1])}{\ensuremath{\beta}}, and design-decay exponent \leanref{sym:kappa}{\ensuremath{\kappa}}, our direct rate is \(n^{-\beta/(2\beta+\kappa+1)}\), as stated in \cref{obj:prop:error-free-reduction}. Definition 2 of \citet{Gaiffas2005} specifies bounded local polynomial-approximation classes, and Theorem 1 characterizes minimax pointwise \(L^p\) risk in direct random-design regression with independent Gaussian response errors. Under its locally symmetric, regularly varying design and regularly varying approximation modulus, the rate combines a power of sample size with a slowly varying factor. Pure-power smoothness and design behavior yield the same exponent as our direct endpoint \citep[Definition 2 and Theorem 1]{Gaiffas2005}. The connection is an exponent comparison across experiments: our target is a population causal mean with bounded outcomes and stratum-specific design envelopes.
% [Gaiffas, Definition 2 and Theorem 1](https://arxiv.org/pdf/math/0410354)

The intermediate inverse regime is related to the regression deconvolution analysis of \citet{FanTruong1993}. Their Sections 3--4 establish estimation rates and corresponding lower bounds for errors-in-variables regression, including pointwise squared-error risk, under positive-design, smoothness, moment, and characteristic-function conditions. In the supersmooth case, exponential attenuation of high frequencies leads to logarithmic estimation rates; Gaussian errors furnish the quadratic-exponential example \citep[Sections 3--4]{FanTruong1993}. Our intermediate branch uses this Fourier inversion mechanism within the weak-design causal class. The comparison therefore concerns how Gaussian attenuation controls localization, with each rate evaluated under its own design restrictions and loss.
% [Fan and Truong, Sections 3--4](https://archive.ymsc.tsinghua.edu.cn/pacm_download/468/10828-Jianqing_Fan_37.pdf)

Compact support supplies another relevant mechanism. \citet[Theorems 1--2]{Meister2007} studies density deconvolution under mean integrated squared error over Sobolev classes, exploiting support information and knowledge of the error characteristic function on a neighborhood of the origin. For source Sobolev smoothness \(s\), polynomial continuation yields the vanishing mean integrated squared error order \(\{(\log\log n)/(\log n)\}^{2s}\). The matching lower bound concerns the joint density-and-error classes of that experiment and carries additional restrictions on density smoothness, the Sobolev radius, and regularity of the local error transform. Our compact-support branch instead addresses pointwise causal evaluation under expected absolute loss, with the Gaussian error law known and the latent design constrained by polynomial envelopes.
% [Meister, Theorems 1--2](https://www.f08.uni-stuttgart.de/.content/media/downloads/Mathematik/mathematische_berichte/2005/2005-007.pdf)

A complementary compact-support result is the spline-assisted regression estimator of \citet{JiangMaCarroll2026}. Under their spline-order, dimension-growth, knot-spacing, bounded-density, compact-support, mean-smoothness, and approximation conditions, together with bounded score-derivative and second-moment operators and the locally unique nonsingular population root in (D5), Theorem 2 gives
\[
\sup_x |\widehat m(x)-m(x)|
=O_p\!\left((nh_b)^{-1/2}+h_b^q\right),
\qquad
E\{\widehat m(x)-m(x)\}^2
=O\!\left(h_b^{2q}+(nh_b)^{-1}\right).
\]
These algebraic spline bounds are conditional on the source experiment's operator and nonsingularity structure \citep[Theorem 2, conditions (C2)--(C4) and (D1)--(D5)]{JiangMaCarroll2026}. In particular, their increasing-dimensional spline system must retain the bounded score operators and locally unique nonsingular population root required by (D5). The measurable polynomial envelopes and Hölder mean restriction used here do not impose that increasing-dimensional operator condition uniformly over the model class. The present result instead gives a same-class minimax characterization of pointwise absolute error and honest interval length, uniform over the maintained Gaussian noise scale; its loss and uniformity quantifiers therefore differ from the source theorem.
% [Jiang, Ma and Carroll, Theorem 2 and its conditions](https://arxiv.org/pdf/1804.00793)

The interval construction connects to two strands of nonparametric inference. \citet[Sections 2--3]{ArmstrongKolesar2020} constructs bias-aware intervals for scalar functionals in direct regression, using a specified smoothness bound to calibrate coverage and assess interval length. Our construction likewise incorporates a uniform bias bound, together with an observable estimation-error certificate. In errors-in-variables regression, \citet[Assumption 3.1 and Theorem 3.2]{KatoSasaki2019} obtains asymptotically valid simultaneous bands through deconvolution and the multiplier bootstrap. Their result assumes ordinary-smooth measurement error, positive design and variance on the reporting interval, moment and regularity conditions, and bandwidth and auxiliary error-sample growth restrictions that include undersmoothing. Alongside these results, \cref{obj:thm:finite-honesty} supplies finite-sample coverage for a connected interval targeting one causal mean, while \cref{obj:thm:uniform-frontier} characterizes its optimal worst-law expected length.
% [Armstrong and Kolesar, Sections 2--3](https://www.princeton.edu/~mkolesar/papers/simple_cis.pdf)
% [Kato and Sasaki, Assumption 3.1 and Theorem 3.2](https://arxiv.org/pdf/1702.03377)

Continuous-treatment causal methods provide the broader identification and estimation context. Propensity-function adjustment identifies population treatment responses under ignorability and positivity \citep[Section 2]{Imai2004}, while doubly robust smoothing combines outcome regression and treatment-density estimation for dose-response estimation \citep[pp.~1229--1230]{Kennedy2017}. Vanishing-noise regression provides a separate precedent for studying rates jointly in sample size and noise scale: \citet[Sections 3--4]{DurotMukherjee2026} analyzes monotone-function estimation in shuffled and unlinked experiments under integrated absolute-error criteria and the respective sampling and error conditions. These connections motivate treating measurement precision as part of the statistical characterization.
% [Imai and van Dyk, Section 2](https://www.ma.imperial.ac.uk/~dvandyk/Research/04-jasa-propfunc.pdf)
% [Kennedy and coauthors, pp. 1229--1230](https://academic.oup.com/jrsssb/article/79/4/1229/7040663)
% [Durot and Mukherjee, Sections 3--4](https://arxiv.org/pdf/2404.09306)

Under the envelope calibration \(c_{\mathrm{lo}}\le(\kappa+1)2^\kappa\le c_{\mathrm{hi}}\), the contribution is a same-class characterization of minimax pointwise expected absolute error and minimum worst-law expected length of honest connected intervals, as defined in \cref{obj:def:minimax-risk,obj:def:honest-length}. For this calibrated class, fixed public \(K\), exponents in their declared ranges, and any noncoverage probability \(\alpha\in(0,1/2)\), \cref{obj:thm:uniform-frontier} gives matching orders uniformly over the maintained Gaussian error range \leanref{sym:sigma}{\ensuremath{\sigma}} in \([0,1/4]\). The observable construction and the finite-sample \((1-\alpha)\)-coverage guarantee in \cref{obj:thm:observable-upper,obj:thm:finite-honesty}, together with the observed-law converse in \cref{obj:thm:observed-lower} under this calibration, place direct estimation, Gaussian Fourier inversion, and compact-support estimation within one causal experiment and one set of model restrictions.

The Gaiffas result supplies a separate direct-regression benchmark for the error-free deduction in \cref{obj:prop:error-free-reduction}; it complements the frontier's own same-class zero-noise conclusion.
